\section{The polygon algebra and its translation groups}\label{sec:geometry}

We first construct the group to which the three main arguments will apply.
Four families of parallel cuts will be used. Two are coordinate cuts; the
other two allow local relations to propagate between widely separated
coordinate cuts.

\subsection{A quadratic ring and four directions}

Put
\[
 \tau=\frac{1+\sqrt5}{2},\qquad \OO=\ZZ[\tau].
\]
The conjugate embedding sends $\tau$ to $\bar\tau=(1-\sqrt5)/2$;
a bar on a vector means conjugation in each coordinate. The unit $\tau$
satisfies $|\bar\tau|=\tau^{-1}$. Choose a positive integer $a$, to be
fixed below, and put $\lambda=\tau^a$. Our four linear forms are
\begin{equation}\label{eq:four-forms}
 \ell_1(x,y)=x,\quad \ell_2(x,y)=y,\quad
 \ell_3(x,y)=y-\lambda x,\quad
 \ell_4(x,y)=x-\lambda y.
\end{equation}
An \emph{allowable line} is $\ell_i^{-1}(t)$ with $t\in\OO$.
Translations by $\OO^2$ preserve these four line families.

\begin{lemma}[Arithmetic of the cuts]\label{lem:arithmetic}
The following statements hold.
\begin{enumerate}[label=\textup{(\roman*)}]
\item There is a positive integer $D$, depending only on $\lambda$, such
that every intersection of nonparallel allowable lines lies in
$D^{-1}\OO^2$. There are finitely many such vertex types modulo
$\OO^2$, including the set of allowable directions through the vertex.
Each type has representatives arbitrarily close to the origin.
\item In each allowable line direction, the group of translations in
$\OO^2$ parallel to that line is dense in the line and has a pair of
generators of arbitrarily small ordinary length.
\item The image of $\OO$ under $z\mapsto(z,\bar z)$ is a lattice in
$\RR^2$. There are absolute constants $C,c>0$ such that, for every
integer $n\ge1$ and every interval $J$ of $n$ consecutive integers,
the points
$\{i\tau\bmod1:i\in J\}$ have largest circular gap at most $C/n$;
any two distinct points in this set have circular distance at least
$c/n$.
\end{enumerate}
\end{lemma}

\begin{proof}
The coefficient determinants of pairs in \eqref{eq:four-forms} belong,
up to sign, to $\{1,\lambda,1-\lambda^2\}$. The first two are units in
$\OO$. Since the reciprocal of a nonzero algebraic integer $v$ is
$\bar v/N(v)$, a common positive integer denominator gives the asserted
$D$. The group $D^{-1}\OO^2/\OO^2$ is finite. Moreover, whether
$\ell_i(v)$ belongs to $\OO$ depends only on this residue class, so
the directions through a vertex introduce no additional infinite
choice. Every coset has dense ordinary image because $\OO^2$ does.

For the coordinate directions the tangent translation groups have
bases $(1,0),(\tau,0)$ and $(0,1),(0,\tau)$. For the other directions
they have bases
\[
 (1,\lambda),\ \tau(1,\lambda)
 \quad\hbox{and}\quad
 (\lambda,1),\ \tau(\lambda,1).
\]
Multiplying both basis vectors by the unit $\tau^{-b}$ leaves the
group unchanged and makes their ordinary lengths as small as desired.

The two vectors $(1,1)$ and $(\tau,\bar\tau)$ generate the stated
lattice, since their determinant is $-\sqrt5$. Multiplication by
$\tau^k$ acts on the two factors with absolute scale factors
$\tau^k$ and $\tau^{-k}$. A fixed lattice covering radius, followed
by such a rescaling, therefore gives a constant $C_0$ with this
property: every axis-parallel rectangle of sufficiently large fixed
area and side lengths comparable to $L^{-1},L$, for $L\ge1$, meets
the lattice. More explicitly, start with a square containing a
translate of a fundamental parallelogram around every point, and
rescale by a unit with $\tau^k$ within a factor $\tau$ of $L$;
increasing one absolute constant absorbs this bounded factor.

Apply this observation with $L$ a small fixed multiple of $n$.
For any chosen ordinary coordinate $t\in[0,1]$, center the long
conjugate interval so that
\[
 i=\frac{z-\bar z}{\tau-\bar\tau}
\]
falls in the middle half of $J$. A conjugate interval of length a
fixed small multiple of $n$, and an ordinary interval of length
$C_1/n$ about $t$, then produce $z=k+i\tau$ with $i\in J$ and
$|z-t|\le C_1/n$. Constants may be enlarged to cover the bounded
small values of $n$. Thus these circular points form a
$C_1/n$-net, which proves the gap bound with $C=2C_1$.

For separation, choose the representative
$z=k+(i-j)\tau$ of a nonzero circular difference with $|z|\le1/2$.
It is a nonzero algebraic integer and
$|\bar z|=|z-(i-j)\sqrt5|\le\sqrt5 n+1/2$.
Consequently
\[
 1\le |N(z)|=|z\bar z|
 \le |z|(\sqrt5 n+1/2).
\]
Taking, for example, $c=(\sqrt5+1)^{-1}$ proves the claim.
\end{proof}

Fix $a$ sufficiently large that
\begin{equation}\label{eq:lambda-choice}
 \lambda c>1000(C+1),\qquad |\bar\lambda|<1/1000.
\end{equation}
This choice is possible because the two embeddings of $\tau^a$ have
reciprocal absolute values. All constants below may depend on this
fixed choice. The two-dimensional version of the lattice statement
is obtained by taking a product: $(z,\bar z)$ for $z\in\OO^2$ is a
lattice in $\RR^4$ with the same simultaneous rescaling property.

\subsection{Boolean polygons and the Stone space}

Consider the Boolean algebra generated by the half-planes bounded by
allowable lines, restricted to a unit square, with opposite sides
identified when applying translations. Two polygonal descriptions are
identified if they differ only on finitely many allowable line
segments. Equivalently, evaluate all descriptions on points avoiding
every allowable line. A nonzero element then has nonempty ordinary
interior. We call its elements \emph{Boolean polygons}; they include
finite unions and complements of ordinary polygonal pieces.

This convention avoids choosing a preferred value at a cut. It can
also be made into an ordinary compact topological space. Let $X$ be
the Stone space of this countable Boolean algebra: its points are
ultrafilters and its clopen sets correspond exactly to Boolean
polygons. The algebra has no atoms, since every region of nonempty
interior can be split by an allowable coordinate cut. Hence $X$ is a
compact metrizable space without isolated points. We continue to
write a polygon $U$ for its associated clopen subset of $X$.
Lebesgue area defines a finitely additive measure $\mu$ on these
sets, with $\mu(X)=1$ and $\mu(U)>0$ whenever $U\ne\varnothing$.

An analogous cut-space realization appears in the Penrose-tiling
full-group construction of
\cite[Section~4 in the author version]{ChornyiJuschenkoNekrashevych2020}:
splitting along line families gives a zero-dimensional space with polygonal
clopen sets.  Here the four line families and their quadratic translation
ring are the ones specified above.

The group
\[
 \Gamma=(\OO/\ZZ)^2\cong\ZZ^2
\]
acts on $X$ by translation modulo one. To see explicitly that this
action is free, use successively finer coordinate rectangles to
associate to every ultrafilter its unique ordinary location in
$\RR^2/\ZZ^2$. This gives a continuous equivariant map from $X$ to
the ordinary torus. A nonzero element of $\Gamma$ has no fixed point
on that torus, and therefore has no fixed point on $X$. The density
of $\Gamma$ on the torus also shows that every orbit meets every
nonempty Boolean polygon: move the ordinary location into its
interior.

All partitions and maps used below have finite descriptions. In
particular, compactness is invoked only to pass from a clopen cover
to a finite clopen partition; it does not permit countably piecewise
group elements.

\subsection{Tracks, translation tables, and slots}

For a positive integer $m$, write $mX=\{1,\ldots,m\}\times X$ and
call its copies of $X$ \emph{tracks}. Extend $\mu$ additively, so
$\mu(mX)=m$. Define $F_m$ to consist of all bijections of $mX$ that
have a finite clopen partition on whose pieces they take the form
\begin{equation}\label{eq:translation-table}
 (i,x)\longmapsto(j,x+\gamma),\qquad \gamma\in\Gamma.
\end{equation}
These maps form a countable group and preserve $\mu$.

Let $U$ be a nonempty Boolean polygon. A \emph{slot with parameter
set $U$} is an embedding $e:U\longrightarrow mX$ given by finitely
many translations and track changes. These slots give the local multisections used to define alternating
full groups in \cite[Section~3]{Nekrashevych2019}.  A collection
$e_1,\ldots,e_r$ of such embeddings with disjoint images realizes
every $\sigma\in\Sym(r)$ as a group element that sends
$e_i(x)$ to $e_{\sigma(i)}(x)$ and fixes the complement. The subgroup
$A_m\le F_m$ is generated by these elements for
$\sigma\in\Alt(r)$, over all finite disjoint slot collections.
It suffices to use three slots and a $3$-cycle, since these generate
each finite alternating group. It also suffices to require each
individual slot to be an ordinary translated copy in a single track:
refine $U$ until all the finitely many embeddings have the form
\eqref{eq:translation-table}, and multiply the resulting permutations
on disjoint pieces.

The allowance of piecewise embeddings makes $A_m$ visibly normal in
$F_m$: conjugation simply composes each slot embedding with an
element of $F_m$. Section~\ref{sec:simplicity} proves that this
alternating subgroup is the desired infinite simple group. The rest
of the paper establishes its two harder properties, amenability and
finite presentation, after a sufficiently large finite $m$ has been
chosen.
